% Transcription — fonds Alexandre Grothendieck, shelfmark 121, batch 1 (pages 1-20).
% Established from the facsimile archives/batches/121/batch-01.pdf (a copy of
% 121.pdf fetched through the project's relay,
% https://grothendieck-relay.onrender.com/source/121.pdf; PDF page k is the
% archivists' page k, no cover sheet), per the skill transcribe-grothendieck.
% The folder is seventy-two pages in four batches (1-20, 21-40, 41-60,
% 61-72), transcribed in parallel; this is the first.
%
% Pass: Opus 5.5 (claude-opus-5-5), 2026-09-26 — first pass, unchecked against the pages by a human.
%
% Numbering. The pencilled numbers bottom-left read 1 on page 1, 4 on page 2
% and then run two ahead (page 9 is pencilled 11, page 20 is pencilled 22):
% pencil 2 and 3 have no scan in this copy. \page{} follows the PDF, as the
% brief fixes; the other batches of the folder should say which they follow.
% He does not paginate this run.
%
% Dating. The folder has its own date (1974, which the inventory title
% attaches to the letter; the notes themselves are « s.d. »); the group rule
% is applied as well, the group « Autour d'Esquisse d'un programme » being
% dated 1968-1991. Nothing on these twenty pages carries a date.
%
% Pages skipped: 2, a blank cream sheet. No letter appears in this batch.
% Pages summarised: none. Page 1 is the grey-blue cover inscribed in ink
% « Topologie Modérée » (his words become the section heading); its lower
% half carries pencil sketches of diagrams and one framed statement, so it
% keeps a \page{1}: the framed statement is transcribed and the sketches are
% described in one \note with what can be read.
% Pages transcribed: 1, 3-20, one blue ink, fast drafting with heavy
% interlinear correction (pp. 5-8, 10, 17 especially).
%
% What the batch is about. A first, set-theoretic approach to tame topology,
% earlier in notation than folder 120 (no R, M_n, (M1)-(M11) here). §1
% (Structure modérée sur un espace): a family F of subsets of a set X,
% stable under finite unions, intersections and complements (a)-c)), the
% algebra A_F of F-step functions with values in Z or a ring Lambda; then X
% topological and the conditions c) (closure; the letter reused), d) (every
% A in F has dense interior in its closure, i.e. boundary rare), e) (no
% infinite strictly decreasing chain of closed members each rare in the
% previous one); Prop: under a)-d)/e) F is generated by its closed members;
% the conditions alpha, beta, gamma, gamma', delta, epsilon on the family F_0
% of closed members, which make sense without the topology, and the coarsest
% topology they define (the closed sets being intersections of members of
% F_0); f) (finitely many connected components, in F, for every locally
% closed member) and local connectedness. §2 (Décompositions cellulaires):
% finite partitions P = (X_i) and F_P; conditions 1) (closures of cells are
% unions of cells) and 2) (distinct cells have distinct closures); cells are
% then locally closed, the incidence order, minimal cells closed; strict
% decompositions (3) connected cells) and equivalence of 3) with f) for F_P;
% the finite quotient space X/R is T_0 (« primitif », doubtful) and X -> X/R
% is open. §3 (F-admissible cellular decompositions: existence): for a finite
% L in F there is a coarsest F-admissible decomposition making every member
% of L a union of cells, proof by induction on the rare set of boundaries;
% corollaries: F is the filtered union of the F_P; sups exist; the variant
% with connected cells under f). §4 (F-regular decompositions): a subfamily
% F_r of « regular » locally closed members, condition (S) (Y_reg dense,
% Y_sing rare), the same existence theorem for F-regular partitions and a
% canonical one. §5 (Localisation): F_A for A in F satisfies a)-e) (and f)
% if F does); the sheaf of germs of members of F, the family F-bar of its
% global sections, F = F-bar when X is quasi-compact and the tame opens form
% a base; example F = {Ø, X}; the question, for F = F_P, whether
% connected cells suffice for F-bar_U = F_U for every tame open U. Ends on
% that question.
%
% Notation fixed for the batch: \mathcal{F}, \mathcal{F}_0 (closed members),
% \mathcal{F}_r, \mathcal{F}_{\mathcal{P}}, \mathcal{F}_A; \mathcal{A}_F for
% his script A; \mathfrak{P}(X) for his power-set P (and \mathfrak{P}_X for
% the sheaf of germs), \mathcal{L} for his script L, \mathcal{P} for the
% partition, \mathcal{R} for the equivalence relation; \complement for his
% complement sign, \dot{A} and bA both for the boundary (he uses both),
% \mathring{A} and A^\circ for the interior; \underline{\mathcal{F}} for the
% sheaf and \overline{\mathcal{F}} for its sections, as he bars them. The
% condition labels are his, including the double use of c) (pp. 3-4) and
% his slips « d) » for f) (pp. 12, 17), noted, not regularised.
%
% Readings to check first: the pencil sketches and framed box of p. 1; the
% heavily reworked condition d) of p. 5; the palimpsest tops of pp. 6, 7, 8
% and 17; the struck block of p. 10 and its left margin; the word read
% « grossière » replacing « fine » on p. 14; « primitif » on p. 13; the
% margins of pp. 3, 8, 10, 11, 15, 18.
% Apparatus: 141 \ill{} and 73 \uncertain{} (counted from the source,
% including those inside \struck{} and \note{}).

\documentclass[11pt,a4paper]{article}
% Le préambule commun à toutes les transcriptions du fonds Grothendieck.
%
% Il tient en une idée : **l'incertitude se compose, elle ne se résout pas.**
% Une transcription qui lisse un mot douteux a détruit la seule chose pour
% laquelle on la faisait — un lecteur doit pouvoir savoir, à chaque endroit,
% ce qui était sur le feuillet et ce qui a été ajouté. Les six macros qui
% suivent sont donc l'appareil critique, et non de la décoration.
%
% Les mêmes commandes sont reconnues par `scripts/render.mjs`, qui produit la
% vue de lecture du site. Une macro ajoutée ici doit l'être aussi là-bas,
% faute de quoi le rendu échouera bruyamment — ce qui est voulu : un rendu
% silencieusement faux est pire qu'un rendu absent.

\ProvidesPackage{grothendieck}[2026/08/08 Transcriptions du fonds Grothendieck]

\usepackage{amsmath,amssymb,amsthm}
\usepackage{mathrsfs}
\usepackage{stmaryrd}
% Les diagrammes commutatifs sont partout dans ce fonds : c'est la langue
% même de la géométrie algébrique de Grothendieck, et une transcription qui
% les rendrait en prose ne transcrirait rien.
\usepackage{tikz-cd}
% Français par défaut — c'est la langue des feuillets — mais l'anglais est
% chargé aussi : la traduction et le résumé sont des documents anglais, et la
% typographie française (espaces avant les deux-points) y serait une faute.
% Ces documents basculent par \selectlanguage{english} en tête de corps.
\usepackage[english,french]{babel}
\usepackage{fontspec}
\usepackage{geometry}
\geometry{a4paper,margin=25mm,marginparwidth=28mm}
\usepackage{marginnote}
\usepackage{xcolor}
% ulem, not soul. `soul`'s \st and \ul are built on a character-by-character
% scan that XeLaTeX's font machinery breaks: the document fails with an
% undefined control sequence rather than degrading. ulem does the same two jobs
% and is Unicode-engine safe. `normalem` stops it hijacking \emph, which the
% transcriptions use for Grothendieck's own emphasis.
\usepackage[normalem]{ulem}
\usepackage{hyperref}
\hypersetup{colorlinks=true,linkcolor=black,urlcolor=black,pdfborder={0 0 0}}

% --- Métadonnées du lot -----------------------------------------------------
% Reprises telles quelles par le script de rendu, qui les affiche en tête de
% la vue de lecture. Elles fixent ce que le fichier prétend couvrir : sans
% elles, une transcription flotte sans point d'attache dans un fonds de seize
% mille feuillets.

\newcommand{\folder}[1]{\gdef\grothFolder{#1}}
\newcommand{\batch}[1]{\gdef\grothBatch{#1}}
\newcommand{\leaves}[2]{\gdef\grothFirst{#1}\gdef\grothLast{#2}}
\newcommand{\foldertitle}[1]{\gdef\grothFoldertitle{#1}}
\gdef\grothFolder{?}\gdef\grothBatch{?}\gdef\grothFirst{?}\gdef\grothLast{?}\gdef\grothFoldertitle{}

% --- Le résumé --------------------------------------------------------------
%
% Il ouvre la lecture modernisée, sous le titre « Résumé », et il est composé
% en petit. Ce n'est pas un ornement : c'est le seul passage du document qui
% ne soit pas des mathématiques, et un lecteur doit pouvoir le reconnaître
% avant de l'avoir lu — puis le sauter s'il connaît déjà le sujet, ou s'y
% arrêter s'il ne le connaît pas. Le filet qui le suit marque la reprise du
% corps.

\newenvironment{resume}
  {\small\setlength{\parskip}{0.45em}}
  {\par\medskip\hrule\medskip}

% --- Mots-clés ---------------------------------------------------------------
%
% En anglais, en clôture du résumé de la lecture modernisée : le vocabulaire
% moderne sous lequel chercher ce que les feuillets construisent. Le manifeste
% du site (`npm run manifest`) les extrait pour étiqueter la cote entière —
% cette ligne est la source unique des tags, il n'y a pas de fichier à côté.
\newcommand{\keywords}[1]{\par\smallskip\noindent
  {\footnotesize\textsc{Keywords} --- \textit{#1}}\par}

% --- L'appareil critique ----------------------------------------------------

% Le numéro de feuillet, en marge. C'est l'ancre qui permet de régler un
% désaccord sur une formule en regardant *un* feuillet plutôt qu'une chemise
% de six cents. Le site s'en sert aussi pour tourner les pages du fac-similé
% quand on fait défiler la transcription.
\newcommand{\leaf}[1]{%
  \marginnote{\footnotesize\color{gray!70}\textsf{#1}}%
  \label{leaf:#1}\ignorespaces}

% Illisible : on ne devine pas. Un mot inventé qui se lit comme les autres est
% le pire résultat possible.
\newcommand{\ill}{\textcolor{red!60!black}{[\,\ldots\,]}}

% Lecture douteuse : le mot est donné, et signalé comme incertain.
\newcommand{\uncertain}[1]{\uline{#1}}

% Ajout de l'éditeur — une lettre manquante, un mot sous-entendu.
\newcommand{\add}[1]{\textcolor{blue!50!black}{[#1]}}

% Biffé par Grothendieck lui-même. Ce qu'il a rayé fait partie du document :
% une rature dit souvent où la pensée a changé de direction.
\newcommand{\struck}[1]{\sout{#1}}

% Note du transcripteur, sur l'état matériel du feuillet.
\newcommand{\note}[1]{\par\smallskip{\small\color{gray!60!black}\textit{[#1]}}\par\smallskip}

% Note marginale de Grothendieck, distincte du corps du texte.
\newcommand{\marginal}[1]{\par\smallskip\noindent\hspace{1em}%
  {\small\color{gray!60!black}#1}\par\smallskip}

% --- Titre ------------------------------------------------------------------

\AtBeginDocument{%
  \begin{center}
    {\large\bfseries Fonds Alexandre Grothendieck --- cote n\textsuperscript{o}~\grothFolder}\\[2pt]
    {\itshape\grothFoldertitle}\\[4pt]
    {\small feuillets \grothFirst--\grothLast\ (lot \grothBatch)}\\[2pt]
    {\footnotesize\color{gray!60!black}Transcription établie d'après le fac-similé numérisé
     par l'Université de Montpellier.\\
     Les lectures incertaines sont soulignées, les illisibles notées [\,\ldots\,],
     les ajouts entre crochets.}
  \end{center}
  \vspace{1em}}

\endinput


\folder{121}
\batch{1}
\pages{1}{20}
\dating{1974}
\watermark{Édition de démonstration}
\foldertitle{Topologie modérée : notes manuscrites (s.d.), lettre (1974)}

\begin{document}

\section*{Topologie modérée}

\page{1}
\note{« Topologie Modérée » est inscrit à l'encre en haut à droite du
feuillet de garde ; ces mots sont de sa main et servent ici de titre. Le bas
du même feuillet porte, au crayon, des esquisses de diagrammes sans texte
suivi, dont on ne donne que ce qui se lit : des cubes $I^{n-1} \to I^n
\leftarrow I^{n+1}$ et des inclusions $X \hookrightarrow I^n$, $Y
\hookrightarrow X$ ; un système $\coprod_{\alpha \in F \subset I} X(\alpha)
\rightrightarrows \coprod_{i \in I} X_i \to X$, avec des flèches $X_i
\xrightarrow{\alpha} X_j$, des $p_i$ vers $X$ et une flèche $X \to Z$ ; une
relation $R \simeq X \times_Y X \subset X \times X$, $R \rightrightarrows X
\to Y$ ; en bas, $|X| \times |Z| \to \Gamma(|g|)$, marqué « cont »,
$(x,z) \mapsto (g(x),z)$, et « $z = \ill{}\, x$ ($= V g(x)$) ».}

\note{au crayon, dans un cadre, en bas à droite du feuillet}
\[
  X \overset{\beta}{\hookrightarrow} M
\]
\begin{itemize}
  \item[a)] $\beta \circ p_i$ modulaires\note{lecture douteuse :
  \uncertain{modulaires} ou \uncertain{modérées} ; suivi de deux mots
  illisibles, \ill{}}
  \item[b)] $\beta$ injectif
\end{itemize}

\page{3}
\section*{§ 1) (Structure modérée sur un espace)}
\note{titre entre parenthèses, souligné d'un trait qui l'enveloppe ; les
lettres a), b), c) des trois conditions sont repassées, et une flèche
courbe les relie}

$X$ ensemble \struck{topologique}

$\mathcal{F}$ ens.\ de parties de $X$, satisfaisant les conditions
\begin{itemize}
  \item[a)] $\mathcal{F}$ stable par réunions finies
  \item[b)] $\mathcal{F}$ stable par intersections finies
  \item[c)] $\mathcal{F}$ stable par passage au complémentaire
\end{itemize}
Si \struck{\ill{}} $\mathcal{A}_{\mathcal{F}}$ est l'ens.\ des fonctions sur
$X$, à valeurs dans $\mathbb{Z}$ (disons), ne prenant qu'un nb fini de
valeurs, et telles que $\forall n \in \mathbb{Z}$, on a $f^{-1}(\lbrace n
\rbrace) \in \mathcal{F}$, \struck{alors de sorte que} alors a) équivaut à
la condition
\begin{itemize}
  \item[a$'$)] $\forall A \in \mathcal{F}$, la fonction caractéristique
  $\varphi_A$ de $A$ est $\in \mathcal{A}_{\mathcal{F}}$ (moyennant a))
\end{itemize}
et par suite la donnée de $\mathcal{F}$ équivaut à celle de la partie
$\mathcal{A}_{\mathcal{F}} \subset \mathbb{Z}^X$. Ceci dit, \uncertain{après}
moyennant a) les conditions b) et b$'$) sont équivalentes, et équivalentes à
la suivante :
\begin{itemize}
  \item[b$''$)] $\mathcal{A}_{\mathcal{F}}$ est stable par produits finis
\end{itemize}
\struck{\ill{} $\mathcal{A}_{\mathcal{F}}$ est \ill{}}\note{une ligne
biffée de hachures ; on n'y distingue que « $\mathcal{A}_{\mathcal{F}}$ est »}
et alors $\mathcal{A}_{\mathcal{F}}$ est \struck{\ill{}} stable par sommes
finies, et \uncertain{s'il l'est} de \uncertain{leurs} $f \circ g$
\ill{} par $f \mapsto -f$, donc c'est un sous-anneau \struck{unitaire} de
$X^{\mathbb{Z}}$\note{sic : on attend $\mathbb{Z}^X$}. Plus généralement,
si $\Lambda$ est un anneau \struck{\ill{}}\add{\uncertain{ou une}
\ill{}}\note{addition interlinéaire peu lisible au-dessus du mot biffé}, on
peut remplacer dans tout ceci \add{$\mathbb{Z}$ par $\Lambda$, donc}
$\mathbb{Z}^X$ par $\Lambda^X$ :
Moyennant a) et b) $\mathcal{A}_{\mathcal{F}}$ est une sous-$\Lambda$-algèbre
de \struck{$\mathcal{A}_{\mathcal{F}}$} $\Lambda^X$. Soit
$\mathrm{Et}(X,\Lambda)$ l'alg.\ des fonctions « étagées » sur $X$ à valeurs
dans $\Lambda$ (i.e.\ ne prenant qu'un \add{nb fini de} \struck{\ill{}}
valeurs). Alors $\mathcal{F}$ est connu quand on connaît la sous-algèbre
$\mathcal{A}_{\mathcal{F}}$ de $\mathrm{Et}(X,\Lambda)$. Lorsque $\Lambda$
est intègre de corps

\marginal{b$'$) \ill{} $\mathcal{A}_{\mathcal{F}}$ \ill{}}\note{note en
biais dans la marge gauche, reliée par une flèche à la condition b$'$), dont
l'énoncé n'est écrit nulle part ailleurs sur la page ; presque entièrement
illisible}

\page{4}
des fractions $K$, on peut prouver que les sous-algèbres $\mathcal{A}$ de
$\mathrm{Et}(X,\Lambda)$ ainsi obtenues sont les traces sur
$\mathrm{Et}(X,\Lambda)$ des sous-algèbres de $\mathrm{Et}(X,\Lambda)
\otimes_\Lambda K \simeq \mathrm{Et}(X,K)$ [i.e.\ satisfaisant la condition
que si $\varphi \in \mathrm{Et}(X,\Lambda)$ \struck{\ill{}} et $\lambda \in
\Lambda$ \struck{\ill{}} \uncertain{sont} tels que $\lambda \neq 0$ et
$\lambda\varphi \in \mathcal{A}$, alors $\varphi \in \mathcal{A}$]

Nous supposons maintenant que $X$ est un espace \emph{topologique}, et
considérons la condition
\begin{itemize}
  \item[c)] $\forall A \in \mathcal{F}$, on a $\overline{A} \in \mathcal{F}$
\end{itemize}
équivalente (moyennant a)) à
\begin{itemize}
  \item[c$'$)] $\forall A \in \mathcal{F}$, on a $\mathring{A} \in
  \mathcal{F}$
\end{itemize}
Ceci implique donc que
\[
  \dot{A} = \overline{A} - A^{\circ} \ \struck{\ill{}} \in \mathcal{F}
\]
donc la \add{\uncertain{prendre}} partition de $X$ en $\mathring{A}$,
$\complement\mathring{A} = X - \overline{A}$, $\dot{A}$ est formée
d'éléments de $\mathcal{F}$\note{la lettre c) de cette nouvelle condition
reprend celle de la condition de passage au complémentaire, p.~3 ; transcrit
tel quel}

\begin{itemize}
  \item[d)] \struck{$\forall A \in \mathcal{F}$, (\ill{} tel que $A$ loc.\
  fermé i.e.\ $\overline{A} - A$ fermé, (i.e.\ $A = \overline{B} - \ill{}$,
  $B, C \in \mathcal{F}$), \ill{} ouvertes, \ill{} fini de) \ill{} les
  composantes connexes sont $\in \mathcal{F}$.}
\end{itemize}
\note{ce d) est barré de trois longs traits obliques ; il sera repris
autrement p.~5}

\page{5}
Condition
\begin{itemize}
  \item[d)] $\forall A \in \mathcal{F}$, \struck{$\mathring{A}$ est dense
  dans $\overline{A}$} \add{l'intérieur de $A$ est dense dans
  $\overline{A}$}, i.e.\ \struck{\ill{}} \add{son complémentaire} \struck{\ill{}
  dans $\overline{A}$, \ill{}} \add{dans $\overline{A}$} est \emph{rare}
  dans $\overline{A}$ (et a fortiori dans $X$)
  \item[e)] \struck{\ill{}} $\forall$ suite décroissante $A_1 \supset A_2
  \supset A_3 \cdots$ de parties fermées \add{$\in \mathcal{F}$}
  \struck{de $X$}, avec $A_{n+1}$ rare dans $A_n$ \uncertain{pour tout} $n$,
  on a $A_n = \emptyset$ pour $n$ assez grand.
\end{itemize}
\note{la première ligne de d) est très raturée : « l'intérieur » est écrit
au-dessus d'une première écriture « l'intérieur de $A$ », et la fin (« est
dense dans $\overline{A}$ », « son complémentaire ») est ajoutée au-dessus de
la ligne ; la reconstitution de l'ordre des ajouts est incertaine}

\emph{NB} Si $\mathcal{L}$ est une partie de $\mathfrak{P}(X)$, alors la plus
petite partie $\mathcal{F}$ de $\mathfrak{P}(X)$ contenant $\mathcal{L}$ et
satisfaisant a) à b)\note{lecture douteuse du renvoi aux conditions}, est
l'ens.\ des réunions finies d'ens.\ de la forme $A \cap \complement B$, avec
$A, B \in \mathcal{L}$. Ceci dit, prouvons

\emph{Prop} Si $\mathcal{F}$ satisfait a) à d), alors $\mathcal{F}$ est
« engendrée » par l'ens.\ des parties fermées $\in \mathcal{F}$.
\note{un trait vertical dans la marge gauche marque l'énoncé}

Soit $\mathcal{F}' \subset \mathcal{F}$ l'ens.\ engendré. \struck{Soit $A
\in \mathcal{F}$, prouvons $A \in \mathcal{F}'$.} Prouvons que \uncertain{si}
$A \in \mathcal{F}$ est $\in \mathcal{F}'$. \struck{Considérons
$\dot{A}$} On a $\overline{A} = \mathring{A} \cup bA$ d'où $A = \mathring{A}
\cup A \cap b\mathring{A}$\note{sic ; on lit $b\mathring{A}$, peut-être pour
$bA$}. $\mathring{A} = \complement\overline{\complement A} \in \mathcal{F}'$
grâce à c) [$\Rightarrow \overline{\complement A} \in \mathcal{F}$ \ill{}
$\overline{\complement A} \in \mathcal{F}'$], il reste à prouver que $A
\cap bA \in \mathcal{F}'$. Mais ceci nous \struck{\ill{}} ramène à prouver
l'énoncé \uncertain{analogue}, avec $A$ remplacé par $bA = X_1$, $\mathcal{F}$
par $\mathcal{F}_1 \subset \mathfrak{P}(X_1)$ « induit » par $\mathcal{F}$.
\add{$A$ par $A_1 = A \cap X_1$} On \uncertain{prend} d), $X_1$ est rare dans
$X$. On construit ainsi une suite décroissante $X_0 = X \supset X_1 \supset
X_2 \supset \cdots \supset X_n \supset \cdots$, avec $X_{n+1}$ rare dans
$X_n$, et $A_n \in \mathcal{F}_n = \mathcal{F} | X_n$, telle que si $A_n
\in \mathcal{F}'_n \Rightarrow A \in \mathcal{F}'$. Or $\exists n$ tel que
$X_n = \emptyset$, \uncertain{cqfd} \ill{}.
\note{« $bA$ » : il écrit $b$ pour la frontière ; plus haut (p.~4) il
notait $\dot{A}$}

\page{6}
Soit $\mathcal{F}_0$ un ens.\ de parties fermées de l'espace top.\ $X$. Pour
que l'ens.\ $\mathcal{F} \subset \mathfrak{P}(X)$ \add{satisfaisant a) b)}
engendré par $\mathcal{F}_0$ soit tel que $\mathcal{F}_0$ soit \add{exactement}
l'ens.\ des parties fermées $\in \mathcal{F}$, il \struck{suffit} \add{faut}
que $\mathcal{F}_0$ soit stable par intersections finies et réunions finies.
Est-ce suffisant, i.e.\ est-il vrai alors que si une réunion finie
\[
  A = \textstyle\bigcup A_i \cap \complement B_i \qquad (A_i, B_i \in
  \mathcal{F}_0)
\]
est \emph{fermée}, alors $A \in \mathcal{F}_0$ ? \struck{\ill{}}\note{la
suite de la ligne, après « $A \in \mathcal{F}_0$ ? », est biffée et
reliée par un trait à la ligne suivante}

On a $A = \overline{A} = \bigcup \overline{A_i \cap \complement B_i}$, donc
il suffit de savoir que $\forall A, B \in \mathcal{F}_0$
\struck{\ill{} réunion \ill{}}, $\overline{A \cap \complement B} \in
\mathcal{F}_0$. \struck{Rédaction \ill{}} \add{\ill{} conditions}
\struck{\ill{}}\note{deux lignes de ratures superposées ; lecture
fragmentaire}

\add{\ill{} donc} que \add{$\mathcal{F}$} \struck{\ill{}} satisfait \ill{},
elle est \add{satisfaite si et seulement si}\note{interligne très chargé ;
lecture incertaine}
\begin{itemize}
  \item[$\alpha$)] $\mathcal{F}_0$ stable par réunions finies
  \item[$\beta$)] --- intersections finies\note{le tiret tient
  lieu de « $\mathcal{F}_0$ stable par »}
  \item[$\gamma$)] $\forall A, B \in \mathcal{F}_0$, on a
  $\overline{A \cap \complement B} \in \mathcal{F}_0$
\end{itemize}
\struck{Montrons que}, \add{\uncertain{alors}} que $\mathcal{F}$ satisfait
d), il faut \struck{que $\mathcal{F}_0$ satisfasse}

\begin{itemize}
  \item[$\delta$)] $\forall A, B \subset \mathcal{F}_0$ \struck{\ill{}
  l'intérieur de} $A \cap \complement B$ dans \ill{}
\end{itemize}
\note{$\delta$) est entouré d'un trait qui le relie à la ligne biffée
au-dessus ; l'énoncé est inachevé}

En effet, si $C = \bigcup A_i \cap \complement B_i$ \add{($C$ \uncertain{
étant} fermé)}, $(A_i, B_i \in \mathcal{F}_0)$, posons $C_i = A_i \cap
\complement B_i$, $D_i = \overline{C_i} \in \mathcal{F}_0$ par $\gamma$),
donc $C_i = D_i - E_i$ ($E_i = D_i - C_i$, fermé rare dans $\mathcal{F}$,
et $\in \mathcal{F}_0$), alors, posant $D = \bigcup D_i$ $= \overline{C}$,
$E_i$ est \struck{rare dans} \add{a fortiori} \struck{\ill{}} rare dans $D
\supset D_i$, donc $E = \bigcup E_i$ est rare dans $D$, donc
\[
  D - E \subset C = \textstyle\bigcup C_i \subset D
\]
\uncertain{donc} $E$ fermé rare dans $D$ donc l'\uncertain{int.} \ill{}.
\add{\uncertain{contient} $D - E$, donc} \ill{} $C$ dans $D = \overline{C}$
\struck{est} \add{\uncertain{est dense dans}} $D = \overline{C}$\ldots

\page{7}
Supposons donc que $\mathcal{F}_0$ satisfasse $\alpha$, $\beta$, $\gamma$.
\struck{Soit $A \in \mathcal{F}$, donc $A = \bigcup C_i$ (\ill{}) $C_i = D_i
- E_i$, $E_i$ rare dans} Alors \struck{$\forall A, B \in \mathcal{F}_0$}
\add{$\forall A \in \mathcal{F}$} \struck{$A \cap \complement B$}
$\overline{A}$ est le plus petit \emph{$\in \mathcal{F}_0$} contenant $A
\cap \complement B$ (donc est connu en termes de $\mathcal{F}_0$ \uncertain{
directement}, sans \uncertain{faire appel} à la top.\ de $X$ !).
\note{les deux premières lignes sont encadrées et biffées ; « $\overline{A}$ »
est écrit au-dessus de « $A \cap \complement B$ » biffé, que la fin de la
phrase conserve}
Donc \struck{la} \struck{\ill{}} condition pour $A \subset B$, $A, B \in
\mathcal{F}$, que $A$ rare dans $B$ i.e.\ $A' = B - A$ dense dans $B$ i.e.\
\struck{\ill{}} $B \subset \overline{A'}$, a un sens indépendamment de la
connaissance de la top.\ de $X$. L'axiome e) sur $\mathcal{F}$ peut donc
s'expliciter en termes de $\mathcal{F}_0$ \uncertain{directement} :
\begin{itemize}
  \item[$\varepsilon$)] $\forall$ suite décroissante $A_1 \supset A_2
  \supset \cdots \supset A_n \supset \cdots$ d'éléments de $\mathcal{F}_0$,
  tels que $\forall n$, $A_{n+1}$ soit \struck{\ill{}} « $\mathcal{F}_0$-rare »
  dans $A_n$ (i.e.\ \struck{$A_n$ est \ill{}} tout $\in \mathcal{F}_0$ qui
  contient $A_n - A_{n+1}$, contient $A_n$), on a $A_n = \emptyset$ pour
  $n$ assez grand.
\end{itemize}
\note{la lettre de cette condition est repassée et peu nette ; on lit
$\varepsilon$), en accord avec la série $\alpha$--$\delta$ de p.~6}

\emph{NB} On peut se demander, indépendamment d'une topologie de $X$, pour
un ensemble $\mathcal{F}_0$ de parties de $X$, quand il existe une top.\
sur $X$ pour laquelle $\mathcal{F}_0$ satisfait : $\alpha$, $\beta$,
$\gamma$) (et \uncertain{éventuellement} $\varepsilon$)). Il faut et il suffit
que $\mathcal{F}_0$ satisfasse $\alpha$, $\beta$) (qui ne font pas appel à
une \emph{topologie} de $X$) et :

\page{8}
\begin{itemize}
  \item[$\gamma'$)] $\forall A, B \in \mathcal{F}_0$, l'ensemble des $\in
  \mathcal{F}_0$ qui contiennent $A \cap \complement B$ est $\in
  \mathcal{F}_0$ (i.e.\ il y a un plus petit $C \in \mathcal{F}_0$ qui
  contienne $A \cap \complement B$).
\end{itemize}
\note{« l'ensemble des $\in \mathcal{F}_0$ » : on attend l'intersection ;
transcrit tel qu'écrit}

C'est suffisant, car on peut \uncertain{alors} prendre la top.\ $T$
\add{dont les fermés sont} \struck{\uncertain{formée}} des intersections
quelconques d'éléments de $\mathcal{F}_0$ [ils sont stables par réunions
finies, comme on voit tout de suite]. Cette topologie \struck{\ill{}}
\uncertain{est} \uncertain{d'ailleurs} la moins fine parmi celles qu'on
envisage. Elle est caractérisée par le fait qu'elle satisfait
\uncertain{à} la condition

\struck{\ill{} les ouverts} \add{(*)} Tout ouvert de $X$ est réunion
d'ouverts de $X$ qui sont $\in \mathcal{F}$, ou encore tout fermé de $X$ est
int.\ de fermés $\in \mathcal{F}$, ou encore \struck{\ill{}} tout $x \in X$
a un syst.\ fond.\ de voisinages $\in \mathcal{F}$ (ou de voisinages
ouverts $\in \mathcal{F}$)\ldots

\emph{NB} Supposons que $\mathcal{F}$ satisfasse a), b), c), d), \uncertain{
i.e.\ que} \struck{\ill{}} si $A \in \mathcal{F}$, alors $\dot{A}$ \add{$=
\overline{A} - \mathring{A}$} est rare dans $X$, \struck{En effet, $\dot{A}
= \overline{A} - \mathring{A}$, \ill{}} \struck{$A^{\circ}$ i.e.\
$\dot{A}$} l'intérieur $U$ de $\dot{A}$ est vide. On sait que $U \in
\mathcal{F}$, nous \uncertain{introduisons} sur $U$ [cela conserve a) b) c)
d)] on est ramené \uncertain{à prouver que} si $\overline{A} = X$,
$\mathring{A} = \emptyset$ \add{$A^{\circ}$ est \ill{} d)} alors \struck{$A$}
$X = \emptyset$ \add{\uncertain{disons donc} $X = \overline{A}$}\ldots
\note{fin de page très raturée et interlignée ; lecture fragmentaire}

\marginal{Condition f) : $\forall A \in \mathcal{F}$, si \ill{} $B, C
\ldots$ \ill{} \uncertain{ferment} (\ill{} $B, C \in \mathcal{F}$) \ill{}
\uncertain{les} \ill{} \uncertain{dans} $A$ \ill{} $\mathcal{F}$}\note{note
écrite verticalement dans la marge gauche, sur trois colonnes, et barrée de
quatre longs traits obliques ; elle est reprise au propre en tête de la
page 9}

\page{9}
\emph{Condition f)} $\forall A \in \mathcal{F}$, $A$ loc.\ fermé\supplied{e} (i.e.\
$A = B - C$, $B, C \in \mathcal{F}_0$) les comp.\ connexes de $A$ sont en
nb fini (donc ouvertes dans $A$) et $\in \mathcal{F}$.
\note{un trait vertical dans la marge gauche marque l'énoncé}

Cela implique donc que $A$ \add{loc.\ fermé $\in \mathcal{F}$} est connexe
ssi \struck{\ill{}} $A', A'' \in \mathcal{F}$, $A' \cap A'' = \emptyset$,
$A' \cup A'' = A$, $A'$ et $A''$ fermés dans $A$ (ou encore dans un sens $A$
d'élément de $\mathcal{F}_0$), on a $A' = \emptyset$ ou $A'' = \emptyset$ :
c'est donc là une condition qui ne fait pas intervenir la top.\ de $X$, mais
$\mathcal{F}_0$ directement. De plus, \struck{$B \subset \ldots$ si $A$ est
\ill{}} \add{\uncertain{$\mathcal{F}_0$ satisfaisant} ($\varphi$)}
\struck{loc.\ \ill{}} si $A = B - C$, avec $B, C \in \mathcal{F}_0$, $C
\subset B$, alors $A$ se décompose en \uncertain{somme} \uncertain{disjointe}
\add{(i.e.\ dans un sens $A$ d'$\in$ de $\mathcal{F}_0$)} \emph{finie}
d'éléments de $\mathcal{F}$ rel.\ fermés dans $A$ \add{qui} sont ainsi
$\mathcal{F}_0$-connexes. \struck{C'est là une \uncertain{nouvelle}
condition (sur $\mathcal{F}_0$).}\note{le ($\varphi$) cerclé, écrit
au-dessus de la ligne, semble nommer la condition qui suit ; lecture
douteuse}

Si les ouverts $\in \mathcal{F}$ forment une base de la top.\ de $X$, alors
la condition f) (qui peut être considérée comme une condition sur l'ens.\ de
parties $\mathcal{F}_0$) implique que les $A \in \mathcal{F}$ (du moins
ceux qui sont loc.\ fermés --- en fait cette restriction est inutile) sont
\emph{loc.\ connexes}.

\page{10}
\section*{§ 2) (Décompositions cellulaires)}
\note{titre entre parenthèses, souligné ; le chiffre 2 est repassé}

(Soit $\mathcal{P} = (X_i)_{i \in I}$ une partition \add{\uncertain{finie}}
de $X$, et soit $\mathcal{F}_{\mathcal{P}}$ l'ens.\ des parties de $X$ qui
sont des réunions \struck{\ill{}} de $X_i$. Alors \struck{\ill{}}
$\mathcal{F}_{\mathcal{P}}$ satisfait les conditions a), b), et
$\mathcal{A}_{(\mathcal{F}_{\mathcal{P}})}$ est l'algèbre des fonctions sur
$X$ à valeurs dans $\Lambda$ qui sont constantes sur chaque $X_i$, i.e.\ qui
se factorisent par l'ens.\ quotient $X/\mathcal{R}$, où $\mathcal{R}$ est la
relation d'équivalence associée à la partition.

\marginal{\emph{NB} La partition \ill{} \uncertain{correspondante} \ill{}
$\mathcal{F}_{\mathcal{P}}$ \ill{} \uncertain{les éléments} \uncertain{minimaux}
\ill{} de $X/\mathcal{R}$ \ill{} \uncertain{réunions} \ill{} $\mathcal{F}_0$
\ill{}}\note{note en biais dans la marge gauche, sur plusieurs lignes,
presque entièrement illisible}

Supposons $X$ un espace topologique. Alors la condition c) équivaut à la
suivante : (1) $\forall i \in I$, $\overline{X_i}$ est une réunion de $X_j$
(i.e.\ est saturé par $\mathcal{R}$).\note{le 1 est cerclé} \struck{Montrons}
\add{\struck{\ill{}} la condition (2) d) équivaut à}
\struck{que cela implique que chaque $X_i$ est \emph{loc.\ fermé} \add{dans
$X$}, i.e.\ ouvert dans $\overline{X_i}$, ou encore que $\overline{X_i} -
X_i$ est fermé. En effet, c'est une réunion de « feuillets » $X_j$, pour $j
\neq i$, et si \add{\uncertain{on a}} ($j \neq i$) \ldots\ $X_j \subset
\overline{X_i}$, \ldots\ dis que l'on \ldots\ $\overline{X_j} \cap X_i =
\emptyset$ : en effet. Je dis que ceci équivaut à}
\note{ce passage est encadré et barré de longs traits obliques}

(2) Si $i, j \in I$, \uncertain{alors} $\overline{X_i} = \overline{X_j}
\Rightarrow X_i = X_j$
\struck{\uncertain{Nécessité} \ill{} Comme c) est satisfaite, chaque $X_i$
est \ill{} loc.\ fermés, donc \ill{}, $X_i = \ill{}$ comme $X_i$ et $X_j$
\ill{} deux \ill{} dans \ill{} (\uncertain{puisque}) $X_i \cap X_j$ est non
vide} \ill{} loc.\ fermé, donc \ill{} dense dans $X_i$. Mais alors
$\overline{X_i} = \overline{X_j} \Rightarrow X_i \cap X_j \neq \emptyset$,
donc $i = j$, OK.
\note{deux lignes et demie barrées ; la fin, sur deux lignes serrées, est
une addition interlinéaire}

Je dis que 1°) et 2°) impliquent que les $X_i$ sont \emph{loc\supplied{alemen}t
fermés}. En effet

\marginal{Mais 1° + les $X_i$ loc.\ fermés implique 2° : $\overline{X_i}
- X_i$ \ill{} $\bigcup_{j} \ill{}$ $(X - \ill{})$ \ill{} $X_i$ \ill{}}
\note{note verticale dans la marge gauche, en bas, peu lisible}

\page{11}
Considérons la relation dans $I$
\[
  i \prec j \quad \overset{\text{déf}}{\Longleftrightarrow} \quad
  \overline{X_i} \subset \overline{X_j}
\]
(relation d'incidence des feuillets). C'est une relation de préordre, et
grâce à 2) une relation d'ordre. Soit \struck{\ill{}} $X_{i_0}$ un feuillet
minimal. Alors, comme $\overline{X_{i_0}}$ donc $\overline{X_{i_0}} -
X_{i_0}$ est une réunion de feuillets, on voit que $\overline{X_{i_0}}
\struck{-} X_{i_0} = \emptyset$ à cause de la condition minimale : \emph{un
feuillet minimal (pour la relation d'incidence) est fermé}. Ceci étant, on
raisonne par récurrence sur le nb de feuillets, en remplaçant $X$ par $X -
X_{i_0}$ et la partition induite $(X_i)_{i \in I - \lbrace i_0 \rbrace}$.
\add{Je \uncertain{dis} maintenant que $\mathcal{F}_{\mathcal{P}}$ est
engendré par un ens.\ de parties fermées $\mathcal{F}_0$, donc d) est
satisfait automatiquement !}
\struck{\ill{}} \emph{Décomposition cellulaire} (finie) de $X$ : satisfaisant
1°) et 2°). Elle est « \emph{stricte} » si elle satisfait de plus
\struck{\ill{}}
\begin{itemize}
  \item[3)] Les $X_i$ sont connexes
\end{itemize}
\note{la phrase « Je dis maintenant \ldots\ automatiquement ! » est serrée
en deux lignes au-dessus de « Décomposition cellulaire », et reliée par une
accolade au paragraphe précédent ; le premier mot de la ligne est
biffé de hachures}

\marginal{\emph{NB} Si $X_i$ \uncertain{est} \ill{} $\overline{X_j}$ \ill{}
dans $X_i \cap \overline{X_j} = \emptyset$, \ill{} $\overline{X_j} \cap
X_i$ \uncertain{rare dans} $X_i$}\note{note verticale dans la marge
gauche, à hauteur du début de la récurrence}

Je dis que [moyennant 1°) et 2°)] 3) équivaut à la condition \textbf{f)} pour
$\mathcal{F}_{\mathcal{P}}$, i.e.\ que \emph{tout} $A \in
\mathcal{F}_{\mathcal{P}}$ qui est loc.\ fermé a \struck{\ill{} loc.\
connexe, et} \add{un nb fini} des comp.\ conn.\ \struck{\ill{}} $\in
\mathcal{F}_{\mathcal{P}}$. \struck{C'est suffisant.} En effet, d) $\Rightarrow$
3), car $X_i$ étant loc\supplied{alemen}t fermé et $\in
\mathcal{F}_{\mathcal{P}}$, \struck{\ill{}}

\page{12}
\struck{\ill{} loc.\ connexe}, et ses composantes connexes \struck{étant}
$\in \mathcal{F}_{\mathcal{P}}$ --- or $X_i$ étant un \add{\uncertain{élément}}
$\neq \emptyset$ minimal de $\mathcal{F}_{\mathcal{P}}$, il s'ensuit que
$X_i$ est connexe. Inversement, 3) $\Rightarrow$ d).\note{sic : d) pour f),
ici et à la ligne précédente}
Soit en effet $A \in \mathcal{F}_{\mathcal{P}}$ \add{\uncertain{loc.\ fermé}}
\struck{(loc.\ fermé)}. Quitte à remplacer $X$ par $A$, et la partition de
$X$ par la partition induite $\lbrace$qui satisfait encore 1), 2), 3)$\rbrace$
--- \struck{c'est} c'est trivial sauf pour 2), \add{mais} \struck{\ill{}
\ill{} $A$ \ill{}}, mais si $X_i, X_j \subset A$, $i \neq j$, et
$\overline{X_i} \cap A = \overline{X_j} \cap A$, \struck{\ill{}} \add{comme}
$X_i \subset \overline{X_i} \cap A$ et $X_j \subset \overline{X_j} \cap A$,
on aura $X_i \subset \overline{X_j}$ et $X_j \subset \overline{X_i}$ donc
$\overline{X_i} = \overline{X_j}$ donc $X_i = X_j$) on est ramené à prouver
\struck{que $X$ est loc.\ connexe et} que les composantes connexes
\struck{\ill{}} de $X$ sont \struck{ouvertes et} réunions de feuillets. Or
les $X_i$ étant connexes, \add{et réunion de $X$,} \struck{les
$\overline{X_i}$ aussi. Donc} les composantes connexes \add{de $X$} sont
\add{bien} des réunions d'ensembles $X_i$, \struck{donc \ill{} 2°), \ill{}
et \ill{} des réunions}.
\struck{Donc il n'y en a qu'un nb fini, donc les composantes connexes sont
\ill{} ouvertes constituant une partition finie de $X$ \ill{} \ill{} des
fermés.}\note{les trois dernières lignes sont barrées de quatre traits
obliques}

Soit $Y = X/\mathcal{R}$ ($\simeq I$) l'espace quotient. C'est un espace top.\
\emph{fini}, \struck{\ill{}} si $i \in I$ $\overline{\lbrace i \rbrace}$ (le
plus petit

\page{13}
ens.\ fermé de $I$ contenant $i$) \add{correspond} \struck{est l'image} au
plus petit ens.\ fermé \emph{saturé} de $X$ contenant $X_i$, i.e.\ à
$\overline{X_i}$ par 1). Donc $j \in \overline{\lbrace i \rbrace}$ ssi $X_j
\subset \overline{X_i}$ i.e.\ $\overline{X_j} \subset \overline{X_i}$ : la
relation d'incidence est la relation de spécialisation. La condition
\add{2°) :} $\overline{X_i} = \overline{X_j} \Rightarrow X_i = X_j$ i.e.\
$i = j$ signifie donc que $I$ est « \uncertain{primitif} »
\struck{\uncertain{condition}} i.e.\ $\overline{\lbrace i \rbrace} =
\overline{\lbrace j \rbrace} \Rightarrow i = j$. Si $\varphi : X \to I$ est
l'application de passage au quotient, la condition 1) signifie que
$\overline{\varphi^{-1}(J)}$ commute à l'adhérence (ou encore, à
l'intérieur) i.e.\ que $\varphi$ est une application \emph{ouverte}.
\note{« primitif » : lecture douteuse (la notion est celle d'un espace
$T_0$) ; « $\overline{\varphi^{-1}(J)}$ commute à l'adhérence » transcrit
tel qu'écrit}
\struck{Donc}, on voit que les décompositions cellulaires \add{(finies)} de
$X$ correspondent aux \struck{\ill{}} espaces quotients \add{finis} de $X$ qui
sont \uncertain{primitifs}, et tels que $X \to X/\mathcal{R}$ soit
\emph{ouvert}.

\section*{§ 3 Décompositions cellulaires $\mathcal{F}$-admissibles :
\uncertain{th.\ d'existence}}
\note{titre ajouté au-dessus de la ligne et entouré ; un premier mot après
« § 3 » est biffé}

Si $\mathcal{F}$ est comme dans § 1, une décomposition cellulaire de $X$
est dite \emph{$\mathcal{F}$-admissible} (ou partition
$\mathcal{F}$-admissible) si ses feuillets sont $\in \mathcal{F}$, i.e.\ si
$\mathcal{F}_{\mathcal{P}} \subset \mathcal{F}$.

\page{14}
\emph{Prop} \struck{\ill{}} Soit $\mathcal{F} \subset \mathfrak{P}(X)$
satisfaisant les conditions a) b) c) \add{d) e)} du § 1, et soit
$\mathcal{L} \subset \mathcal{F}$ une partie finie de $\mathcal{F}$. Alors il
existe une décomposition cellulaire $\mathcal{F}$-admissible de $X$, telle que
les $Y \in \mathcal{L}$ sont des réunions de feuillets. Parmi toutes ces
décompositions cellulaires, il y en a une plus \struck{fine}
\add{grossière} (\struck{\ill{}} \add{qui est donc unique déterminée}).
\note{le mot remplaçant « fine » est peu lisible ; « grossière » est offert
d'après le sens}

\marginal{\uncertain{OK} \ill{} \uncertain{(fini)} \ill{} \ill{}}\note{dans
l'angle supérieur gauche, quelques mots en biais, barrés de hachures}

\struck{\ill{}} Soit $X' = \bigcup_{Y \in \mathcal{L}} \dot{Y}$, c'est une
réunion finie de parties fermées rares de $X$, donc est une partie rare de
$X$, soit \struck{$X' = X_{\ill{}}$} $U = \complement X'$. Les traces
\add{$Y \cap U$} sur $U$ des $Y \in \mathcal{L}$ sont des \struck{ouvertes}
parties à la fois ouvertes et fermées \struck{\ill{}} en \emph{nb fini}. Il y
a donc (\add{ensemblistement} \uncertain{immédiat}) dans $U$ une plus
\add{grossière} \struck{\ill{}} partition \add{(finie)} par des $U_i$
\struck{\ill{}} telles que les $Y \cap U$ soient des réunions \add{(finies)}
de $U_i$ ; de plus, chaque $U_i$ est une \uncertain{intersection} \add{(finie)}
de \uncertain{tels} ensembles $Y \cap U$, donc est aussi ouvert et fermé.
Ceci étant, si la prop.\ est prouvée pour \struck{\ill{}} $(X, \mathcal{L})$ remplacé par $(X', \mathcal{L}' = \lbrace Y
\cap X' \mid Y \in \mathcal{L} \rbrace \cup \lbrace \overline{U_i} \cap X'
\rbrace_{i \in I})$\note{la seconde moitié de $\mathcal{L}'$ est écrite
sous la ligne ; la barre sur $U_i$ est peu nette},


\page{15}
il est immédiat qu'il l'est pour $(X, \mathcal{L})$ : si $(X'_j)_{j \in J}$
est la partition cherchée de $X'$, celle de $X$ est formée des $X'_j$ et des
$U_i$. Ceci nous fournit un principe de dém.\ par récurrence, si on suppose
que $\mathcal{F}$ satisfait la condition e)
\struck{(Dim) Pour toute suite \add{finie} $X_1 \supset X_2 \supset \cdots
\supset X_n$ de parties fermées $\in \mathcal{F}$ de $X$, telles que
$X_{i+1}$ soit rare dans $X_i$, on a $X_n = \emptyset$ pour $n$ assez
grand.}\note{ce paragraphe, étiqueté « (Dim) » dans la marge, est barré de
trois traits obliques ; c'est la condition e) de la p.~5}
\marginal{\uncertain{résulte} \uncertain{des} \ill{} de (Dim) \ill{}}\note{en
biais dans la marge gauche, à côté du paragraphe barré}

\emph{Cor 1} Si $\mathcal{F}$, \add{satisfait les} \struck{en plus des
conditions} a) b) c) \add{d) e)} \struck{, satisfait la condition (Dim)},
alors $\mathcal{F}$ est réunion filtrante des sous-ens.\ \add{de la forme}
$\mathcal{F}_{\mathcal{P}}$ (associés aux décompositions cellulaires
\struck{\ill{}} $\mathcal{F}$-admissibles de $X$).

\emph{Cor 2} L'ens.\ \add{$\Sigma$} des décompositions cellulaires
$\mathcal{F}$-admissibles, pour la relation évidente de raffinement, admet des
bornes sup.\ pour les \supplied{sous-}ens.\ finis de \add{$\Sigma$}
\struck{\ill{}}.

\page{16}
\emph{Remarque} Supposons que $\mathcal{F}$, en plus de a) b) c) \add{d) e)}
\struck{(Dim)} satisfasse aussi f). Alors il existe une décomposition
cellulaire \emph{stricte} $\mathcal{P}$ (i.e.\ à feuillets connexes)
$\mathcal{F}$-admissible telle que les $Y \in \mathcal{L}$ soient des
\add{réunions de} feuillets, et parmi ces $\mathcal{P}$, il y en a une plus
grossière. La dém.\ est la même, mais au lieu de prendre les $U_i$
précédents, on prend leurs composantes connexes. Donc $\mathcal{F}$ est
réunion filtrante des $\mathcal{F}_{\mathcal{P}}$, avec les $\mathcal{P}$
décompositions cellulaires à \emph{feuillets connexes}, et ces
décompositions cellulaires admettent des sup finis.

\section*{§ 4 (Décompositions cellulaires $\mathcal{F}$-modérées
\uncertain{régulières})}
\note{titre écrit au-dessus de la ligne, entre parenthèses et souligné}

(Supposons maintenant donnée une partie $\mathcal{F}_r$ de $\mathcal{F}$
formée d'éléments de $\mathcal{F}$ \add{loc.} \struck{\ill{}} fermés : on
les appelle les parties \emph{$\mathcal{F}$-régulières} de $\mathcal{F}$. On
suppose que
\begin{itemize}
  \item[(S)] Pour tout $Y \in \mathcal{F}$, $Y$ \struck{\ill{}} fermé, parmi
  les parties \add{rel.} ouvertes $U$ de $Y$ qui sont $\in \mathcal{F}_r$, il
  y en a une plus grande, $Y_{\mathrm{reg}}$, et celle-ci est dense dans $Y$
  i.e.\ son complémentaire $Y - Y_{\mathrm{reg}} = Y_{\mathrm{sing}}$ est rare
  dans $Y$.
\end{itemize}
\note{la lettre de la condition est cerclée et repassée ; on lit (S), que la p.~17 confirme (« la condition (S) sur $\mathcal{F}_r$ »)}

\page{17}
Ceci dit, une décomposition cellulaire $\mathcal{P}$ de $X$ est dite
\emph{$\mathcal{F}$-régulière} si \add{$\mathcal{F}_{\mathcal{P}}$ \ill{}}
\struck{les feuillets} \add{\ill{} conditions de régularité induites par
$\mathcal{F}$, $\mathcal{F}_r$} \struck{sont $\in \mathcal{F}_r$ \ill{}}
\add{satisfait la condition S, et} \struck{\ill{}} \add{\uncertain{\ill{}
précédemment}}, \struck{\ill{}} $\forall Y$ fermé $\mathcal{P}$-saturé
$Y_s$ est $\mathcal{P}$-saturé. (Cela implique que les feuillets sont $\in
\mathcal{F}_r$.)
\note{les trois premières lignes sont un palimpseste de ratures et
d'additions interlinéaires ; lecture très partielle}

\emph{Prop} Sous les conditions a) b) c) d) \add{e)} \struck{\ill{}} sur
$\mathcal{F}$, et la condition (S) sur $\mathcal{F}_r \subset \mathcal{F}$,
si $\mathcal{L} \subset \mathcal{F}$ est une partie finie, il existe une
partition $\mathcal{F}$-\emph{régulière} de $X$, et il y en a une plus
\emph{grossière} telle partition. Si $\mathcal{F}$ satisfait \emph{en plus}
f), alors il existe une partition $\mathcal{F}$-\emph{régulière} de $X$ à
feuillets connexes, et une plus \emph{grossière} telle.
\note{un trait vertical dans la marge gauche marque l'énoncé ; « telle que
les $Y \in \mathcal{L}$ soient des réunions de feuillets » n'est pas écrit}

OPS $\mathcal{L}$ stable par $\complement$.

Soit $X' = \bigcup_{Y \in \mathcal{L}} \dot{Y} \cup Y_{\mathrm{sing}}$, $U = X
- X'$. Les $Y \cap U$ sont ouverts et fermés, et réguliers. On continue comme
plus haut, \add{avec les $U_i$, mais dans $\mathcal{L}'$ on met aussi les
$(\overline{U_i})_{\mathrm{sing}}$} dans le cas où $\mathcal{F}$ satisfait
d), on prend \add{avant} aux composantes connexes des $U$.
\note{« d) » pour f), sic}

\emph{Cor 1} Donc \add{$\mathcal{F}$} est réunion filtrante de
$\mathcal{F}_{\mathcal{P}}$, où $\mathcal{P}$ est à feuillets
$\mathcal{F}$-réguliers (resp.\ $\mathcal{F}$-réguliers et connexes, si
$\mathcal{F}$ satisfait d)).

\emph{Cor 2} Sup.\ \add{finis} dans les partitions $\mathcal{F}$-régulières
(resp.\ à fibres connexes)

\emph{Cor 3} Partition $\mathcal{F}$-régulière canonique de $X$

\page{18}
\struck{§ 5 \ill{}}\note{une ligne de titre entièrement biffée de
zigzags}
\section*{§ 5 Localisation}

Soit $\mathcal{F}$ satisfaisant a) b) c) d) e) ou encore $\mathcal{F}_0$
satisfaisant $\alpha\,\beta\,\gamma\,\varepsilon$ (c'est kif-kif). Soit $A \in
\mathcal{F}$, considérons $\mathcal{F}_A = \lbrace B \in \mathfrak{P}(A) \mid
B \in \mathcal{F} \rbrace$. Évidemment $\mathcal{F}_A$ satisfait a) b) c),
satisfait-il aussi d) et e) ? Pour d), c'est \uncertain{prouvé} au fond
\uncertain{puisqu'il} est immédiat que $\mathcal{F}_A$ est engendré par
$(\mathcal{F}_A)_{\mathrm{f}}$, i.e.\ que tout $\in \mathcal{F}_A$ est
réunion finie de parties loc.\ fermées \add{$\in \mathcal{F}_A$} de $A$. \ill{}
Enfin, pour e) on considère $A \supset A_1 \supset A_2 \supset \cdots \supset
A_i \supset \cdots$ $A_i \in \mathcal{F}_A$, $A_i$ rel\supplied{ativemen}t fermé
dans $A$, $A_{i+1}$ rare dans $A_i$. Alors $\overline{A_i}$ rare dans
$\overline{A_{i+1}}$ \struck{\ill{} pour $i$ grand} (car $A_i - A_{i+1}$
dense dans $A_i$ donc dans $\overline{A_i}$, \struck{\ill{} $A_{i+1}$} or
$A_i - A_{i+1} = A_i - \overline{A_{i+1}} \cap A_i \subset \overline{A_i} -
\overline{A_{i+1}}$, qui est donc dense dans $\overline{A_i}$) donc pour $i$
grand on a $\overline{A_i} = \emptyset$ donc $A_i = \emptyset$ cqfd.
\note{« $\overline{A_i}$ rare dans $\overline{A_{i+1}}$ » : sic, pour
$\overline{A_{i+1}}$ rare dans $\overline{A_i}$ ; sous
« $\overline{A_{i+1}} \cap A_i$ » est écrit « $A_{i+1}$ »}

\marginal{On peut \uncertain{toujours} \ill{} \uncertain{restreindre} à une
\ill{} \uncertain{quelconque} $A$ de $X$, \uncertain{en posant}
$\mathcal{F}_A = \lbrace B \ldots \rbrace$ \ill{} $(B \in \mathcal{F})$, dans
$\mathcal{F}_A$ \uncertain{satisfait} a) b) c) d) e)}\note{en biais dans la
marge gauche ; lecture très incertaine}

Prouvons que si $\mathcal{F}$ satisfait aussi f), il en est de même de
$\mathcal{F}_A$. C'est trivial si $A$ loc.\ fermé. Il suffit de prouver que
si $A \in \mathcal{F}$ quelconque, ses comp.\ conn.\ sont $\in \mathcal{F}$ et
en nb fini. Or $A = \bigcup Z_i$, les $Z_i$ loc.\ fermés, et par f) OPS les
$Z_i$ connexes. Donc les composantes

\page{19}
connexes de $A$ sont réunions de $Z_i$, donc en nb fini et $\in
\mathcal{F}$.

Nous supposons maintenant que les ouverts $\in \mathcal{F}$ forment une base
de la top.\ de $X$. Considérons chaque partie de $X$ comme une section du
faisceau \add{$\mathfrak{P}_X$} des germes de parties de $X$. Le
sous-faisceau \add{$\underline{\mathcal{F}}$} \struck{\ill{}} de
$\mathfrak{P}_X$ engendré par $\mathcal{F}$ (considéré comme ens.\ de
sections) \add{admet} \struck{comme sections} \add{comme sections sur $X$
les} $A \subset X$ satisfaisant la condition
\begin{itemize}
  \item[a)] $\forall x \in X$, $\exists$ voisinage ouvert $U$ de $x$, tel que
  $U \in \mathcal{F}$ et $A \cap U \in \mathcal{F}$, ou encore : $X$ est
  réunion d'ouverts $U_i$ tels que $U_i \in \mathcal{F}$ et $U_i \cap A \in
  \mathcal{F}$.
\end{itemize}
\note{un trait vertical dans la marge gauche marque la condition ; la
lettre a) est repassée}

Soit $\overline{\mathcal{F}}$ l'ens.\ de ces $A$. On a $\mathcal{F} \subset
\overline{\mathcal{F}}$, $\overline{\overline{\mathcal{F}}} =
\overline{\mathcal{F}}$. On voit aussi $\overline{\mathcal{F}}$ satisfait a)
b) c) d). \struck{(satisfait-il e) ?)} Il n'est pas clair qu'il satisfait
e), \struck{\ill{}} ni que $\overline{\mathcal{F}}$ est engendré par les
fermés qu'il contient. Mais si $X$ est \emph{quasi-compact}, on a
$\mathcal{F} = \overline{\mathcal{F}}$ (car on peut prendre les $U_i$ en nb
fini).
\note{le d) est suivi d'un indice ou d'une virgule peu nette}

Si on ne suppose plus que la top.\ de $X$ soit engendrée par les ouverts
$\in \mathcal{F}$, on

\page{20}
peut néanmoins définir $\underline{\mathcal{F}}$, ses sections sur l'ouvert
$U \subset X$ sont les $A \subset U$ tels que $\forall x \in U$, $\exists$
voisinage (ouvert si on veut) $V_x$ et $B_x \in \mathcal{F}$ tels que $A \cap
V_x = B_x \cap V_x$. Posant encore $\overline{\mathcal{F}} = \Gamma(X,
\underline{\mathcal{F}})$ (plus généralement, $\Gamma(U,
\underline{\mathcal{F}}) = \overline{\mathcal{F}}_U$, si $U$ ouvert de $X$),
satisfait a) b) c) d) (peut-être pas e), ni $\overline{\mathcal{F}}$
engendré par ses fermés\ldots), mais, si $X$ compact, on ne peut affirmer que
$\overline{\mathcal{F}} = \mathcal{F}$. Par exemple, si $\mathcal{F} =
\lbrace \emptyset, X \rbrace$, alors $\overline{\mathcal{F}}$ est formé des
parties à la fois ouvertes et fermées de $X$ (\struck{donc} si $X$ est loc.\
connexe ce sont les réunions de comp.\ connexes).
\note{« si $X$ compact, on ne peut affirmer » : sic ; on attendrait « même
si »}

\struck{Plus généralement}

Supposons p.\ ex.\ que $\mathcal{F} = \mathcal{F}_{\mathcal{P}}$, où
$\mathcal{P}$ est une décomposition cellulaire de $X$. Quand a-t-on
$\overline{\mathcal{F}} = \mathcal{F}$ ? Toute partie ouverte et fermée
d'une cellule fermée est $\in \overline{\mathcal{F}}$, il faut donc que les
cellules fermées soient connexes. Plus généralement pour \emph{tout} ouvert
$U \subset X$, $U \in \mathcal{F}$, $\overline{\mathcal{F}}_U =
\mathcal{F}_U$, on voit donc qu'il faut que \emph{toutes} les cellules
soient connexes (puisqu'elles sont fermées dans un ouvert $U \in
\mathcal{F}$). Cela suffit-il pour que $\overline{\mathcal{F}_U} =
\mathcal{F}_U$ pour tout $U$ ouvert $\in \mathcal{F}$ ?

\end{document}
